\section{Results}

\subsection{Visualization Results}

The three different data sets show different levels of signal detection.


In the first Cas A data set (Figure~\ref{fig: CasA1_Data}, there is a peak in trace 1, but it is not differentiated enough from trace 2 to confidenty call it a hydrogen emission.


In the second Cas A data set from April 7 (Figure~\ref{fig: CasA_040726_Data}, there is a strong peak in trace 1 that is not present in trace 2. This indicates that there is a potential hydrogen emission found during the observation.


In the Virgo A data set (Figure~\ref{fig: VirgoA Data}, a sharp spike appears in trace 2, but not in trace 1. This is most likely caused by some form of interference or noise. 


\subsection{Analysis Results}

After averaging and smoothing the two traces, that data became significantly less noisy, making it easier to identify the peaks in the frequency spectra.
The smoothing process was able to reduce random fluctuations and highlight the signal.


For the first Cas A data set (Figure~\ref{fig:casa_analysis}), there was still a peak present, but it was much less distinct. The signal-to-noise ratio was lower, which indicated that the signal was weaker and closer to noise.
It's possible this was a detection, but it is much less certain compared to the second Cas A data set.


For the second Cas A data set from April 7 (Figure~\ref{fig:casa0407_analysis}), the smoothed data showed a clear peak that stood above the noise level. This resulted in a high signal-to-noise ratio, which confirms that there was a detected signal that could represent hydrogen emission.



For the Virgo A data set (Figure~\ref{fig:virgoa_analysis}), smoothing did not really reveal a consistent signal. Instead, sharp spikes remained that appeared in trace 1 and 2/
Once again, these features are most likely the cause of some form of interference, and there was no reliable detection. 
